\chapter{Guía de Despliegue y Manual de Usuario}
\label{anx:despliegue}

\section{Requisitos previos}

\section{Configuración de credenciales}

\section{Arranque con Docker Compose}

\section{Manual de usuario}

\chapter{Detalles Técnicos}
\label{anx:tecnico}

\section{Arquitectura de agentes}

\subsection{Jerarquía de clases y registro}
\label{anx:jerarquia-agentes}

Todos los agentes de \textbf{Seshat} heredan de \texttt{\_BaseAgent}, que encapsula la lógica
de reintentos, el límite de concurrencia y la invocación estructurada del LLM. Sobre esta base
se especializan dos familias:

\begin{itemize}
  \item \textbf{Agentes de identificación} (\texttt{\_BaseIdentificationAgent}): extienden
        \texttt{\_BaseAgent} con un esquema de salida genérico parametrizado por el tipo de
        concepto. Cada agente concreto (\texttt{DecisionIdentificationAgent},
        \texttt{RiskIdentificationAgent}, etc.) sólo necesita declarar su
        \texttt{ConceptType}, el modelo de salida Pydantic y el prompt de sistema. Añadir un
        nuevo tipo de concepto se reduce a crear un fichero con estas tres declaraciones y
        registrarlo en el \texttt{AgentRegistry}.

  \item \textbf{Agentes de resolución} (\texttt{BaseSameTypeResolutionAgent},
        \texttt{BaseCrossTypeResolutionAgent}): siguen el mismo patrón. El agente de mismo
        tipo parametriza los tipos de relación válidos para su \texttt{ConceptType}; el de tipo
        cruzado declara el par \texttt{(source\_type, target\_type)} que gestiona. El
        \texttt{ResolutionRegistry} instancia en paralelo los nueve agentes de tipo cruzado y
        uno por tipo para la resolución de mismo tipo.
\end{itemize}

El \texttt{AgentRegistry} descubre los agentes disponibles en tiempo de ejecución a partir de
\texttt{ExtractionConfig.concept\_types}, lo que permite habilitar o deshabilitar tipos de
concepto mediante configuración sin modificar el código del orquestador.

\subsection{Matriz de relaciones por par de tipos}
\label{anx:matriz-relaciones}

La tabla~\ref{tab:rel-matrix} recoge las relaciones que un agente de resolución puede asignar
para cada par dirigido (tipo fuente~$\rightarrow$~tipo destino). La diagonal corresponde a la
resolución de mismo tipo; el resto, a la de tipo cruzado. Un guión (\texttt{---}) indica que
ningún agente gestiona ese par en la versión actual: de las 12 combinaciones dirigidas entre
tipos distintos, sólo 9 se habilitan.

\begin{table}[H]
  \centering
  \small
  \begin{tabularx}{\textwidth}{l*{4}{>{\raggedright\arraybackslash}X}}
    \toprule
    \textbf{Fuente $\downarrow$ / Destino $\rightarrow$}
      & \texttt{decision} & \texttt{risk} & \texttt{action\_item} & \texttt{open\_question} \\
    \midrule
    \texttt{decision}
      & \texttt{supersedes}, \texttt{amends}, \texttt{conflicts\_with}
      & \texttt{mitigates} & \texttt{blocks} & \texttt{resolves} \\
    \addlinespace
    \texttt{risk}
      & \texttt{blocks} & \texttt{amends} & \texttt{blocks} & \texttt{blocks} \\
    \addlinespace
    \texttt{action\_item}
      & ---
      & \texttt{mitigates}
      & \texttt{supersedes}, \texttt{amends}, \texttt{conflicts\_with}, \texttt{blocks}, \texttt{depends\_on}
      & --- \\
    \addlinespace
    \texttt{open\_question}
      & \texttt{blocks} & --- & \texttt{blocks} & \texttt{amends}, \texttt{depends\_on} \\
    \bottomrule
  \end{tabularx}
  \caption{Relaciones permitidas por par de tipos (fuente~$\rightarrow$~destino). La diagonal
           es resolución de mismo tipo; el resto, de tipo cruzado.}
  \label{tab:rel-matrix}
\end{table}

\subsection{Técnicas de diseño de prompts}
\label{anx:tecnicas-prompts}

El diseño de los prompts de los agentes de identificación combina varias técnicas
complementarias:

\begin{itemize}
  \item \textbf{Razonamiento en cadena} (\textit{chain-of-thought}~\cite{wei2022cot}): el
        prompt descompone la tarea en pasos intermedios verificables —localizar el fragmento
        relevante, evaluar las compuertas lógicas, clasificar el tipo— antes de emitir la
        salida final. Hacer explícito el razonamiento reduce los errores de clasificación al
        exponer los pasos intermedios a la propia lógica del modelo.

  \item \textbf{Anclaje en cita literal} (\textit{quote-first grounding}): antes de describir
        un nodo, el agente debe localizar y copiar el fragmento literal del texto fuente que lo
        respalda. Esto reduce la alucinación al obligar al modelo a anclar cada afirmación en
        evidencia concreta antes de sintetizarla.

  \item \textbf{Compuertas lógicas}: el prompt enumera condiciones necesarias y suficientes
        para emitir un nodo —por ejemplo, para que un \texttt{risk} sea válido deben cumplirse
        simultáneamente cuatro condiciones: modo de fallo concreto, implicación sustantiva del
        grupo, no totalmente resuelto en el mismo intercambio, y no un requisito sin consecuencia
        declarada. Si alguna condición no se satisface, el nodo no debe emitirse.

  \item \textbf{Contraejemplos etiquetados}: el prompt incluye ejemplos de casos límite que
        \emph{no} deben extraerse, con explicación de por qué. Esto delimita la frontera
        semántica entre categorías próximas (por ejemplo, \texttt{risk} frente a
        \texttt{open\_question}).

  \item \textbf{Delimitadores XML}: el contenido de la transcripción y los \textit{hints} de la
        base de conocimiento se encierran en bloques \texttt{<transcript>} y
        \texttt{<kb\_hint>}. Las instrucciones del sistema indican explícitamente que cualquier
        texto con apariencia de instrucción dentro de esos bloques debe tratarse como datos y
        no ejecutarse. Esto mitiga la inyección de prompts desde el contenido de la reunión.

  \item \textbf{Salida estructurada forzada}: los esquemas Pydantic se inyectan como
        instrucción de herramienta en la llamada al modelo. Las violaciones de tipo activan
        reintentos automáticos, lo que elimina la ambigüedad de formato y hace el
        post-procesamiento determinista.
\end{itemize}

\subsection{Agentes reflectivos}
\label{anx:agentes-reflectivos}

Ambas familias de agentes soportan un modo de \textbf{auto-revisión} (\textit{self-review})
opcional que añade una segunda pasada de LLM sobre la salida del agente primario.

El \texttt{ReflectiveIdentificationAgent} implementa un ciclo \emph{extraer → validar →
filtrar}: tras la extracción, una llamada de validación comprueba el cumplimiento lógico y
semántico de cada nodo (¿satisface las reglas de extracción? ¿coincide la descripción con la
cita?); los nodos que fallan se descartan. En caso de fallo de la validación, todos los nodos
se devuelven sin filtrar, garantizando que el modo reflectivo nunca sea más perjudicial que el
modo plano.

El \texttt{ReflectiveResolutionAgent} actúa como árbitro para los casos genuinamente ambiguos.
Sólo las entradas con \texttt{alt\_rel\_type} no nulo se envían a una llamada de desempate;
las entradas claras se omiten por completo, lo que minimiza el coste adicional. Ante cualquier
fallo del árbitro, se conserva el \texttt{rel\_type} original.

Ambos agentes utilizan el patrón \textit{subclass-and-delegate}: heredan de la misma base que
los agentes planos y delegan todas las propiedades abstractas al agente interno. El orquestador
y los registros de agentes no requieren ningún cambio para soportar el modo reflectivo.

\subsection{Caché de prompts y optimización de latencia}
\label{anx:cache-prompts}

Los prompts de sistema de los agentes son estáticos por tipo de concepto y se reutilizan en
todos los trabajos, lo que los hace idóneos para caché de prompts a nivel de proveedor. En
Anthropic, las cabeceras \texttt{cache\_control} marcan el prompt de sistema como candidato a
caché; en OpenAI, el prefijo compartido se cachea automáticamente. La transcripción, en cambio,
no se cachea de forma deliberada: los agentes de identificación se ejecutan en paralelo, de modo
que todas las llamadas son concurrentes y fallarían la caché, pagando la prima de escritura sin
ninguna lectura posterior. La efectividad de la caché
—\textit{cache\_read\_input\_tokens} y \textit{cache\_creation\_input\_tokens} por llamada— se
registra en MLflow como métricas observables. El primer trabajo tras arrancar el worker incurre
en coste de creación de caché (\textit{cold start}); los trabajos posteriores se benefician de
la reducción de latencia y coste.

\section{Capa de observabilidad}
\label{anx:observabilidad}

La medición del consumo y la latencia descrita en la sección~\ref{sec:observabilidad} se apoya en
dos rastreadores —\texttt{UsageTracker} y \texttt{LatencyTracker}— que se instalan como
\textit{callbacks} de LangChain. El reto de implementación es que los agentes se ejecutan de forma
concurrente: la solución evita propagar el rastreador por la firma de cada método y lo publica en
una variable de contexto (\texttt{contextvars.ContextVar}). Las tareas hijas heredan el valor de la
variable al crearse, de modo que todos los agentes lanzados dentro de una etapa acumulan en el
mismo rastreador sin cambio alguno en sus firmas.

Cada etapa del \textit{pipeline} se envuelve con dos decoradores —\texttt{track\_token\_budget} y
\texttt{track\_latency\_profile}— que, al entrar, crean el rastreador y lo fijan en la variable de
contexto y, al salir, registran los totales como métricas de MLflow (tokens por tipo, y latencia
mínima, media, mediana, p95 y máxima). El decorador de presupuesto lee sus límites de la
configuración en tiempo de llamada, avisa al alcanzar el 90\,\% del cupo y aborta la etapa al
superarlo en más de un 10\,\%; este margen tolera que varios agentes concurrentes rebasen el cupo
antes de que cualquiera de ellos pueda observarlo. Los totales de cada etapa se agregan además en
un rastreador a nivel de trabajo.

Dos fuentes de consumo no pasan por las llamadas al LLM y se contabilizan mediante envoltorios
transparentes: \texttt{TrackingEmbeddings} envuelve el modelo de \textit{embeddings} y cuenta sus
tokens de entrada, y \texttt{TrackingTranscriber} envuelve al transcriptor y contabiliza los
segundos de audio procesados. Ambos empujan su medición al rastreador activo en la variable de
contexto, integrándose en el mismo presupuesto que las llamadas al LLM.

\section{Arquitectura de evaluación}

\subsection{Estructura común de los arneses y el emparejador de identificación}
\label{anx:harnesses-eval}

Los cinco arneses comparten el mismo esqueleto de ejecución, condicionado por un detalle del marco
de evaluación: \texttt{mlflow.genai.evaluate} invoca su función de predicción de forma síncrona,
mientras que los componentes de \textbf{Seshat} —orquestador y agentes— son asíncronos. Para
salvar esa incompatibilidad, cada arnés pre-computa todas las predicciones antes de entregar el
control a MLflow: un método \texttt{\_run\_all\_predictions} recorre el corpus, ejecuta el
componente real de cada caso y almacena el resultado en una caché en memoria; la función de
predicción que ve MLflow es entonces una mera consulta a esa caché. Sobre las predicciones
cacheadas, MLflow aplica el \textit{scorer} del arnés —una función determinista, sin llamadas al
LLM—, agrega las métricas por caso y el arnés vuelca el veredicto al fichero de \textit{gate}
mediante \texttt{upsert\_gate}, que actualiza sólo el bloque del arnés en ejecución y conserva el
resto.

El emparejador de identificación merece una descripción aparte, pues es la pieza no trivial de esa
cadena: debe decidir qué nodo extraído corresponde a qué nodo esperado antes de poder contar
aciertos y fallos. Resuelve un emparejamiento bipartito \textit{greedy} restringido a pares del
mismo tipo de concepto. Para cada par candidato calcula una puntuación de solapamiento en $[0,1]$
por una de dos vías: si el nodo predicho lleva cita (\texttt{quote\_anchors}), la puntuación es la
similitud difusa (\texttt{rapidfuzz}, \textit{partial ratio}) entre la cita esperada y el fragmento
de transcripción que abarca la predicha; si no la lleva, se recurre a una similitud difusa
ponderada sobre título y descripción. Los pares se ordenan de mayor a menor puntuación y se van
reclamando de forma \textit{greedy} —cada nodo esperado y cada predicho se asignan una sola vez—;
un par sólo se acepta como acierto si supera el umbral \texttt{QUOTE\_MATCH\_THRESHOLD} de 0.70. Los
nodos predichos sin pareja cuentan como falsos positivos y los esperados sin pareja, como falsos
negativos. Sobre los pares aceptados, el \textit{scorer} puntúa además la exactitud de los campos
estructurados propios de cada tipo (difusa para \texttt{assignee}, \texttt{due}, \texttt{rationale}
y \texttt{context}; exacta para el subtipo de \texttt{risk}; solapamiento de conjuntos para
\texttt{alternatives\_considered}), señales que se registran en MLflow pero no condicionan el
\textit{gate}.

Cada arnés registra en la traza de MLflow, junto a la salida que consume el \textit{scorer}, una
vista de comparación legible: los nodos —o relaciones— esperados y predichos uno al lado del otro.
El objetivo es la depurabilidad: en la interfaz de MLflow cada caso es una fila, y un caso fallido
debe poder diagnosticarse desde esa fila sin abrir el fichero YAML del corpus. Un
post-procesador de trazas recorta además los campos más verbosos de los nodos —citas en bruto,
descripciones largas— y conserva sólo tipo, título, descripción y confianza, de modo que la
comparación quepa de un vistazo.

\subsection{Caché de evaluación y meta-scorers de calibración}
\label{anx:cache-eval}

La caché de predicciones que sostiene la ejecución de los arneses es también la que permite
calibrar umbrales sin repetir llamadas al LLM. Su pieza central es \texttt{read\_or\_run}: dada la
ruta de un fichero de caché y la corrutina que produciría el resultado, devuelve el contenido
cacheado si el fichero existe —cerrando la corrutina sin ejecutarla— y, en caso contrario, ejecuta
la corrutina y persiste el resultado como JSON. El nombre del fichero se deriva de tres
componentes: el identificador del caso de corpus, una huella (\textit{fingerprint}) del agente y un
\textit{hash} de la entrada. La huella del agente incorpora su \textit{prompt} y configuración, de
modo que cualquier cambio en el agente invalida automáticamente sus entradas de caché sin afectar a
las de otras variantes. Al terminar, cada arnés no borra la caché, sino que aplica un barrido de
marcado (\texttt{sweep\_stale\_entries}) que elimina únicamente las entradas en ámbito que no se
tocaron en la ejecución —predicciones obsoletas por un cambio de \textit{prompt} o de entrada—,
preservando las válidas para ejecuciones posteriores.

Sobre esa caché operan los dos meta-\textit{scorers} de calibración. Ambos siguen un diseño en dos
fases: primero pueblan la caché con el resultado del \textit{pipeline} para cada caso —reutilizando
las predicciones que el arnés correspondiente ya hubiera generado—, y después barren el umbral
sobre esos resultados en memoria, sin más entrada/salida ni llamadas al modelo.

El meta-\textit{scorer} de identificación barre el umbral sobre la confianza heurística en el
intervalo $[0,1]$ con paso 0.01. En cada punto reproduce qué nodos se auto-aprobarían y, aplicando
el emparejador de identificación contra la verdad de referencia, calcula la precisión entre los
aprobados y la cobertura, de forma agregada y por tipo de concepto. El umbral sugerido es el que
maximiza la cobertura sujeto a que la precisión alcance un objetivo configurable (0.95 por
defecto); si ninguno lo logra, se conserva el de mayor precisión. Los empates se resuelven hacia el
umbral más bajo. Deliberadamente no se emplea F1: un nodo por debajo del umbral no es un error, sólo
se difiere a revisión humana.

El meta-\textit{scorer} de \textit{retrieval} barre el umbral de similitud con paso 0.005. Para que
un único resultado cacheado sirva a cualquier corte, las búsquedas se siembran sin filtro de
puntuación y el umbral se aplica en memoria durante el barrido. El umbral sugerido es el que
maximiza el F2 macro-promediado: los casos positivos aportan su F2 —que pondera el \textit{recall}
el doble que la precisión— y los negativos aportan especificidad (1 si no se devuelve nada por
encima del umbral, 0 en caso contrario), lo que impide que un umbral en cero, que lo aceptaría
todo, domine la selección.

\section{Referencia de la API REST}
\label{anx:api-reference}

Todos los \textit{endpoints} cuelgan del prefijo de versión \texttt{/v1}. La
tabla~\ref{tab:api-endpoints} los agrupa por router. Salvo el router \texttt{/admin} —autenticado
con la clave raíz— y el sondeo de salud, todos exigen una clave de \ac{API} válida y el rol indicado
en la sección~\ref{sec:orquestacion} y en los mecanismos de seguridad.

\begin{longtable}{@{}l l >{\raggedright\arraybackslash}p{0.42\textwidth}@{}}
  \caption{\textit{Endpoints} de la API REST, agrupados por router (prefijo \texttt{/v1}).}
  \label{tab:api-endpoints} \\
  \toprule
  \textbf{Método} & \textbf{Ruta} & \textbf{Descripción} \\
  \midrule
  \endfirsthead
  \toprule
  \textbf{Método} & \textbf{Ruta} & \textbf{Descripción} \\
  \midrule
  \endhead
  \midrule
  \multicolumn{3}{r}{\footnotesize\textit{continúa en la página siguiente}} \\
  \endfoot
  \bottomrule
  \endlastfoot

  \multicolumn{3}{@{}l}{\textbf{\texttt{/jobs} — ciclo de vida del \textit{pipeline}}} \\
  \addlinespace
  \texttt{POST}   & \texttt{/jobs}                     & Crear un \textit{pipeline} a partir de un archivo de entrada \\
  \texttt{GET}    & \texttt{/jobs}                     & Listar \textit{pipelines} con filtros \\
  \texttt{GET}    & \texttt{/jobs/\{id\}}              & Consultar el estado de un \textit{pipeline} \\
  \texttt{GET}    & \texttt{/jobs/\{id\}/results}      & Recuperar los nodos extraídos \\
  \texttt{GET}    & \texttt{/jobs/\{id\}/transcript/excerpt} & Extraer un fragmento de la transcripción \\
  \texttt{POST}   & \texttt{/jobs/\{id\}/approve}      & Enviar decisiones de revisión y disparar la escritura \\
  \texttt{POST}   & \texttt{/jobs/\{id\}/retry}        & Reintentar un \textit{pipeline} fallido \\
  \addlinespace
  \multicolumn{3}{@{}l}{\textbf{\texttt{/graph} — consulta de la base de conocimiento}} \\
  \addlinespace
  \texttt{GET}    & \texttt{/graph}                    & Consultar nodos por filtros \\
  \texttt{GET}    & \texttt{/graph/search}             & Búsqueda semántica sobre los nodos \\
  \texttt{GET}    & \texttt{/graph/\{id\}}             & Obtener un nodo por su identificador \\
  \texttt{GET}    & \texttt{/graph/\{id\}/neighbours}  & Vecinos directos (profundidad~1) \\
  \texttt{GET}    & \texttt{/graph/\{id\}/detail}      & Nodo junto con sus vecinos directos \\
  \texttt{GET}    & \texttt{/graph/\{id\}/impact}      & Travesía multi-salto de impacto \\
  \texttt{POST}   & \texttt{/graph/nodes}              & Añadir un nodo manualmente \\
  \texttt{POST}   & \texttt{/graph/nodes/bulk}         & Añadir nodos en lote \\
  \texttt{POST}   & \texttt{/graph/nodes/resolve}      & Resolver relaciones para nodos dados \\
  \texttt{PUT}    & \texttt{/graph/nodes/\{id\}}       & Editar un nodo \\
  \texttt{PUT}    & \texttt{/graph/nodes/\{id\}/override} & Anular contenido con justificación \\
  \texttt{DELETE} & \texttt{/graph/nodes/\{id\}}       & Eliminar un nodo \\
  \texttt{DELETE} & \texttt{/graph/nodes/bulk}         & Eliminar nodos en lote \\
  \texttt{GET}    & \texttt{/graph/relationships}      & Listar relaciones \\
  \texttt{POST}   & \texttt{/graph/relationships}      & Crear una relación \\
  \texttt{DELETE} & \texttt{/graph/relationships/\{id\}} & Eliminar una relación \\
  \addlinespace
  \multicolumn{3}{@{}l}{\textbf{\texttt{/admin} — provisión de claves (clave raíz)}} \\
  \addlinespace
  \texttt{GET}    & \texttt{/admin/api-keys}           & Listar claves de \ac{API} \\
  \texttt{POST}   & \texttt{/admin/api-keys}           & Crear una clave de \ac{API} \\
  \texttt{DELETE} & \texttt{/admin/api-keys/\{id\}}    & Revocar una clave de \ac{API} \\
  \addlinespace
  \multicolumn{3}{@{}l}{\textbf{Identidad y salud}} \\
  \addlinespace
  \texttt{GET}    & \texttt{/me}                       & Resolver la clave de \ac{API} a usuario y rol \\
  \texttt{GET}    & \texttt{/health}                   & Sondeo de salud de la \ac{API} \\
  \texttt{GET}    & \texttt{/health/components}        & Estado de las dependencias externas \\
\end{longtable}

\section{Esquema de la base de datos}

\section{Prompts clave}

\section{Configuración completa de referencia}
